% ASSERT_CONVENTION: natural_units=dimensionless, metric_signature=mostly_minus,
%   jordan_product=(1/2)(ab+ba), octonion_basis=fano_e1e2=e4,
%   complex_structure=u_equals_e7, potential=minus_log_det,
%   cone_metric=hess(minus_log_det), cubic_norm_det=ring_lemma_verification_det_3,
%   curvature_engine=totaro_closed_form, arithmetic=exact_over_Q
%
% Phase 72 (B -- Matter-Sourcing), Plan 01 (DERV-02 + VALD-04): the matter-aware
%   Totaro curvature engine on the dim-4 h_2(C_u) spacetime sub-slice and the M=0
%   baseline, for the v17.0 milestone "Gravity as Intrinsic Curvature of the h_3(O)
%   Bulk Geometry".
%
% STATUS (2026-05-31): 72-01 STOPPED AT THE FIRST-RESULT GATE; the milestone has
%   BACKTRACKED to Phase 70 (researcher decision, see Sec on adjudication). The
%   plan's required first load-bearing result -- "the M=0 center is Einstein
%   (traceless Ricci S_munu = 0) with R(center) = -3 < 0 (pure-Lambda Cartan)" --
%   DOES NOT HOLD, and this is the plan's explicit backtracking trigger
%   (disconfirming observation #3). The finding is decisive, exact-over-Q, and
%   cross-checked against the engine's own H^3 benchmark: the 4-dim cone-Hessian
%   sub-slice M=0 vacuum is a static product R x H^3 whose FLAT factor (1,1,0,0) is
%   the TIMELIKE x_0 direction under h_2(C_u) ~ R^{3,1}; its Ricci endomorphism has
%   eigenvalues {0,-1,-1,-1}, so it is NOT Einstein in n=4.
%
%   THIS NON-EINSTEIN M=0 VACUUM *IS* THE VALD-04 RESULT (negative-result-is-success),
%   NOT a reference-choice problem. A GR Lambda-vacuum must be Einstein (Ric prop. to
%   g); no baseline subtraction turns a non-Einstein vacuum into an Einstein one, so
%   "which Lambda reference" presupposes a maximally-symmetric vacuum this
%   construction does not produce. The matter Delta R(M) verdict deliverables are NOT
%   produced (honest-negative discipline; no masking SUMMARY). The matter mechanism
%   and the Delta-R isolation pipeline are PRESERVED (they are not what broke) and
%   stand for whichever metric the reopened Phase 70 selects.
%
% References:
%   [TOT04] Totaro, "The curvature of a Hessian metric," arXiv:math/0401381.
%   [FK94]  Faraut & Koranyi, Analysis on Symmetric Cones (OUP 1994), Ch. II-IV.
%   [McC04] McCrimmon, A Taste of Jordan Algebras (Springer 2004).
%   [71VER] .gpd/phases/71-a-homogeneity-kill-gate/71-VERIFICATION.md (matterless
%           regression anchors; cone-Hessian R(center) = -3; K_sect = -1/2).
%   [70-02] derivations/70-signature-bridge.tex + code/bulk_geometry_verification.py:
%           construction-(ii) bridge; Hess(-log det)|_{I/3} = diag(9,9,18,18)/26244.

\section{Phase 72-01: matter mechanism, Phase-71 regression, and the M=0 non-Einstein vacuum (VALD-04 finding $\Rightarrow$ backtrack to Phase 70)}

\subsection{Conventions (carried verbatim from the v17.0 lock)}

All decisive quantities are EXACT over $\mathbb{Q}$ (no float on any verdict; ranks
via \texttt{sympy.Matrix}). The cone metric is the Hessian of the potential
$\Phi=-\log\det_3$, with $\det_3$ the certified Freudenthal cubic norm (engine SSOT,
cross-term order $2\,\mathrm{Re}((x_2 x_1)x_3)$; \texttt{octonion\_algebra.py}
banned). The Riemann tensor is the Totaro closed form
\begin{equation}
  R_{ijkl} = -\tfrac14\, g^{pq}\big(C_{jlp}C_{ikq}-C_{ilp}C_{jkq}\big),
  \qquad C_{ijk}=\Phi_{,ijk},
  \label{eq:totaro}
\end{equation}
with $\mathrm{Ric}_{jl}=g^{ik}R_{ijkl}$, $R=g^{jl}\mathrm{Ric}_{jl}$. The spacetime
sub-slice is $h_2(C_u)$, engine-native indices $\{1,2,3,10\}=(\beta,\gamma,p,q)$,
mostly-minus. The center is $X_{\mathrm{bg}}=I/3$.

\subsection{Peirce index map and the unique matter channel (Task 1, CLEAN)}
\label{sec:mechanism}

Under $E_{11}=\mathrm{diag}(1,0,0)$ the multiplication operator $L_{E_{11}}$ is
diagonal in the engine-native basis, with Peirce eigenspaces (verified exactly over
$\mathbb{Q}$ via the $L_{E_{11}}$ spectrum)
\begin{equation}
  V_1=\{0\}\ (\alpha),\qquad
  V_0=\{1,\dots,10\}=h_2(\mathbb{O}),\qquad
  V_{1/2}=\{11,\dots,26\},
\end{equation}
matching the module constants \texttt{SPACETIME\_SLICE\_IDX}$=[1,2,3,10]$,
\texttt{V\_HALF\_IDX}$=[11..26]$, \texttt{V1\_ALPHA\_IDX}$=[0]$. In the engine layout
$x_1=X_{21}$ (the $V_0$ octonion, components $\{3..10\}$), $x_2=X_{02}$ ($\{11..18\}$),
$x_3=X_{10}$ ($\{19..26\}$). The cubic norm is
\begin{equation}
  \det_3(X)=\alpha\beta\gamma-\alpha|x_1|^2-\beta|x_2|^2-\gamma|x_3|^2
            +2\,\mathrm{Re}\big((x_2 x_1)x_3\big).
  \label{eq:det3}
\end{equation}
Subtracting the diagonal-norm part of~\eqref{eq:det3} leaves \emph{exactly}
$2\,\mathrm{Re}((x_2x_1)x_3)$ (verified as a symbolic identity over $\mathbb{Q}$ on
the full 27-coordinate $X$). The diagonal $\alpha,\beta,\gamma$ -- including the
$V_1$ direction $\alpha=x_0$ -- do NOT appear inside the octonion triple; they couple
only to the squared norms $|x_1|^2,|x_2|^2,|x_3|^2$. Hence:
\begin{itemize}
  \item the triple $2\,\mathrm{Re}((x_2x_1)x_3)$ is the \emph{unique} $V_0\!\leftrightarrow\!V_{1/2}$
        octonion-mixing channel (this is what "cross-terms off" means precisely;
        the $|x_2|^2,|x_3|^2$ norm couplings to $\alpha,\beta,\gamma$ remain --
        Open Question 2, stated, not over-claimed);
  \item $\alpha$ ($V_1$) is absent from the triple, structurally explaining the
        Phase-71 finding that $V_1$ is inert to the slice curvature.
\end{itemize}
The symmetric trilinear polarization of the cubic norm,
$d(D_{x_2},D_{x_1},D_{x_3})$, on Fano-connected octonionic basis directions
$(x_1{:}e_4,\,x_2{:}e_5,\,x_3{:}e_7)$ evaluates to $2\neq0$ over $\mathbb{Q}$ (and to
$0$ when any one slot is zeroed), confirming the cubic vertex coupling matter to
$V_0$ is genuinely present and not vacuously zero.

\subsection{Non-vacuity preflight: the chosen representative $M$ (Task 1, CLEAN)}
\label{sec:nonvacuity}

The representative non-vacuous matter placement carried to 72-02 is
\begin{align}
  \texttt{bg\_delta} &= \{\,4{:}\tfrac14,\ 7{:}\tfrac15\,\}
    &&\text{(FIXED $V_0$-internal $x_1$: $e_1,e_4$ content),}\\
  \texttt{matter\_delta} &= \{\,11{:}\tfrac13,\ 15{:}\tfrac16,\ 19{:}\tfrac13,\ 23{:}\tfrac17\,\}
    &&\text{($M$ in $V_{1/2}$: $x_2,x_3$ each with $e_0,e_4$).}
\end{align}
At $X_{\mathrm{bg}}+M$ all three octonion slots $x_1,x_2,x_3$ are populated, each with
$e_4$ content (genuinely octonionic / non-associative), and the cross-term is nonzero:
\begin{equation}
  2\,\mathrm{Re}\big((x_2 x_1)x_3\big)\Big|_{X_{\mathrm{bg}}+M} = -\tfrac{13}{315}\neq 0
  \quad\text{(exact over $\mathbb{Q}$).}
\end{equation}
This defeats the hollow-triple-product pitfall (Pitfall 3): the cross-term off-switch
of 72-02 will be a genuine control.

\subsection{Phase-71 matterless regression anchors (Task 1, CLEAN)}
\label{sec:regression}

Using the verified matterless path \texttt{\_curvature\_invariants\_at(delta, slice\_pt)}
with zero matter ($\delta$ in $V_0$-internal $\{4..9\}$ only) and the common fixed
slice point \texttt{ROUTE1\_COMMON\_SLICE\_PT}$=[\tfrac25,\tfrac35,\tfrac1{10},\tfrac18]$,
the Phase-71 [71VER] Ricci-scalar anchors are reproduced EXACTLY over $\mathbb{Q}$:
\begin{center}
\begin{tabular}{lll}
\hline
$\delta$ (matterless) & $R$ (Ricci scalar) & status\\
\hline
$\{4{:}\tfrac13\}$        & $-73041507/21967969$    & matches [71VER] A(ii)\\
$\{4{:}\tfrac15,5{:}\tfrac17\}$ & $-521269105/154700283$  & matches [71VER] A(i)\\
$\{4{:}\tfrac25,6{:}\tfrac13\}$ & $-137053539/48874081$   & matches [71VER] A(i)\\
$\{7{:}\tfrac13,8{:}\tfrac14,9{:}\tfrac15\}$ & $-114049215/37982569$ & matches [71VER] A(i)\\
\hline
\end{tabular}
\end{center}
and at the center ($\delta=\{\}$, center slice $[\tfrac13,\tfrac13,0,0]$),
$R(\text{center})=-3$, $\det g = 26244$. The matter-aware path of 72-01 is therefore a
faithful extension of the verified matterless engine (regression passes).

\subsection{The M=0 baseline obstruction: the cone-Hessian sub-slice is a metric
            cone over $H^3$, NOT Einstein (FIRST-RESULT GATE)}
\label{sec:obstruction}

The plan's required first load-bearing result is that the $M=0$ center be Einstein,
$S_{\mu\nu}\equiv R_{\mu\nu}-\tfrac{R}{4}g_{\mu\nu}=0$, with $R=-3<0$ (pure-$\Lambda$
Cartan). Computing the Totaro curvature~\eqref{eq:totaro} of the cone-Hessian metric
$g_{ij}=\mathrm{Hess}(-\log\det_3)$ restricted to the dim-4 slice at the center
$X_{\mathrm{bg}}=I/3$ (exact over $\mathbb{Q}$) gives $R=-3$ as required, but the
Ricci tensor is NOT proportional to the metric:
\begin{equation}
  g = \mathrm{diag}(9,9,18,18),\qquad
  \mathrm{Ric} =
  \begin{pmatrix} -\tfrac92 & \tfrac92 & 0 & 0\\
                  \tfrac92 & -\tfrac92 & 0 & 0\\
                  0 & 0 & -18 & 0\\
                  0 & 0 & 0 & -18\end{pmatrix},\qquad
  g^{-1}\mathrm{Ric}\ \text{has eigenvalues}\ \{0,-1,-1,-1\}.
  \label{eq:ricci-center}
\end{equation}
The Ricci endomorphism $g^{-1}\mathrm{Ric}$ has the degenerate spectrum
$\{0,-1,-1,-1\}$ (so $S_{\mu\nu}\neq0$; the slice is NOT Einstein in $n=4$). The
eigenvalue-$0$ eigenvector is $(1,1,0,0)$, the $\beta+\gamma$ direction. On the
Riemannian cone this is the \emph{dilation / overall-scale} direction -- exactly the
scale-invariance $R(X)=R(2X)$ found in Phase 71 ($\Phi=-\log\det$ shifts by a constant
under $X\mapsto cX$, leaving the Hessian metric, hence $R$, unchanged). \emph{But under
the Lorentzian spacetime reading} $h_2(C_u)\simeq\mathbb{R}^{3,1}$ with the $\det_2$
Minkowski form $\beta\gamma/3-p^2/3-q^2/3$, the very same $(1,1,0,0)=\beta{+}\gamma$
direction is the \textbf{timelike $x_0$ direction} (the $\beta\gamma/3$ block is the
hyperbolic $1{+}1$ part; $(1,1)$ is its positive/timelike eigenvector). So the flat
factor of the cone is \emph{time}, and the $M=0$ vacuum is a static product
$\mathbb{R}_{\text{time}}\times H^3_{\text{space}}$. The coordinate sectional
curvatures confirm the structure: $K(e_0,e_1)=0$ (the timelike/dilation 2-plane is
flat), the mixed planes are $-\tfrac14$, and $K(e_2,e_3)=K(p,q)=-\tfrac12$ -- the
$H^3$ base value, identical to the engine's certified
\texttt{h3\_cone\_hessian\_benchmark}.

\paragraph{Diagnosis (decisive, exact over $\mathbb{Q}$, not an engine bug).} The
4-dim cone-Hessian sub-slice is a \emph{metric cone} $C(H^3)=\mathbb{R}_{>0}\times H^3$:
a flat radial (dilation) line times the constant-curvature $H^3$ base. A non-trivial
metric cone in dimension $n=4$ is never Einstein -- its Ricci endomorphism has a
radial $0$-eigenvalue and base eigenvalues shifted, here $\{0,-1,-1,-1\}$ -- so the
\emph{conjunction} the contract demands ("$R=-3$ AND $S_{\mu\nu}=0$") cannot hold for
this metric. The two candidate baseline metrics each satisfy only one half:
\begin{itemize}
  \item the cone-Hessian $g_{ij}=\mathrm{Hess}(-\log\det_3)$ (the metric Phase-71's
        $R(\text{center})=-3$ anchor refers to) has $R=-3$ but is a cone, hence NOT
        Einstein ($S_{\mu\nu}\neq0$, eigenvalues $\{0,-1,-1,-1\}$);
  \item the construction-(ii) inherited metric $g_{\mu\nu}=\eta_{\mu\nu}+h_{\mu\nu}$
        reduces at the center to the \emph{constant} Minkowski form $g=\eta_{\mathrm{bg}}$
        ($h=0$ at the center by construction), which is FLAT ($R=0$), trivially
        Einstein but with $\Lambda=0$, contradicting $R=-3$.
\end{itemize}
Per the disconfirming-observation discipline (this is named obstruction #3) and the
first-result-gate rule, 72-01 STOPS here: the matter $\Delta R(M)$ verdict is not
produced, and no result is reported that papers over the non-Einstein vacuum. The
clean Task-1 results (Secs~\ref{sec:mechanism}--\ref{sec:regression}) stand. The
adjudication (Sec~\ref{sec:adjudication}) records that this non-Einstein $M=0$ vacuum
\emph{is} the VALD-04 finding and that the milestone backtracks to a reopened Phase 70
(Sec~\ref{sec:phase70-task}) to decide which object is the physical spacetime metric.

\subsection{Adjudication: the non-Einstein vacuum IS the result; the candidate
            ``baselines'' are not interchangeable (researcher decision, 2026-05-31)}
\label{sec:adjudication}

The obstruction is \textbf{not} a reference choice. A general-relativistic
$\Lambda$-vacuum \emph{must} be Einstein ($\mathrm{Ric}\propto g$); no baseline
subtraction turns a non-Einstein vacuum into an Einstein one. So the question ``which
$\Lambda$ reference?'' already presupposes a maximally-symmetric vacuum that this
construction \emph{does not produce}. The computed $M=0$ cone-Hessian vacuum is the
static product $\mathbb{R}_{\text{time}}\times H^3$ with Ricci spectrum
$\{0,-1,-1,-1\}$ -- and \textbf{that is the VALD-04 finding}
(negative-result-is-success), to be reported as such, not papered over.

The three ``reconciliations'' previously listed are not on equal footing:
\begin{itemize}
  \item[\textbf{(a)}] \textbf{REJECTED -- ``Einstein on the $\det_2=1$ $H^3$ leaf.''}
        This references 4-dim gravity against the 3-dim \emph{spatial} $\det_2=1$
        hyperboloid, i.e.\ it \emph{drops the flat $=$ timelike $x_0$ direction}. $H^3$
        is maximally symmetric but it is the \emph{wrong} symmetric space and one
        dimension short: GR's empty hyperbolic-slice vacuum is Milne $=$ flat
        Minkowski, \emph{not} $H^3$. Measuring matter-sourcing against it manufactures
        a pass by swapping spacetime for a spatial slice -- this is the
        \texttt{fp-relabel} class of self-deception. Do not take it. (Quotienting the
        dilation/timelike direction -- old option 3 -- is the same move and is
        likewise rejected.)
  \item[\textbf{(b)}] \textbf{CONDITIONAL -- ``$g=\eta+h$ is the decisive metric.''}
        $\eta+h$ is flat at the center, which \emph{is} GR's correct empty vacuum
        (right physics, unlike (a)). But it is admissible only with two caveats:
        (i) the flatness is \emph{inserted} by the center-subtraction, so $\Lambda=0$
        is built in, not derived -- Phase 73's per-equation circularity audit must
        carry this tripwire; and (ii) adopting $\eta+h$ as the spacetime metric
        \emph{abandons} the route thesis ``gravity $=$ curvature of the cone-Hessian,''
        since the cone-Hessian gives $\mathbb{R}\times H^3$, not flat. Hence (b) is
        admissible ONLY after a reopened Phase 70 explicitly rules that $\eta+h$ (not
        the cone-Hessian) is THE spacetime metric and restates the thesis accordingly.
\end{itemize}

\textbf{DECISION (c): halt the matter-sourcing measurement and reopen Phase 70.} The
plan's backtracking trigger \#3 has fired; it is honored, not routed around. The
matter mechanism (Sec~\ref{sec:mechanism}), the non-vacuity preflight
(Sec~\ref{sec:nonvacuity}), and the $\Delta R$-isolation pipeline
(Sec~\ref{sec:carryforward}) are PRESERVED and are not what broke; they stand for
whichever metric Phase 70 selects.

\subsection{Reopened Phase 70: which object is the physical spacetime metric?}
\label{sec:phase70-task}

The Riemannian cone-Hessian (the thesis metric $\Rightarrow$ a
$\mathbb{R}\times H^3$, non-Einstein $M=0$ vacuum) and the Lorentzian construction-(ii)
bridge $\eta+h$ ($\Rightarrow$ a flat $M=0$ vacuum) \emph{disagree about the $M=0$
geometry}. That disagreement is precisely the construction-(ii) signature bridge being
\emph{underdetermined} -- the Riemannian/Lorentzian tension already flagged in
\texttt{CONVENTIONS.md}. The reopened Phase 70 must resolve exactly one question:
\emph{which object is the physical spacetime metric?}
\begin{itemize}
  \item \textbf{If the cone-Hessian (the thesis):} the $M=0$ vacuum is non-Einstein
        ($\mathbb{R}\times H^3$), so the route does NOT yield GR-with-a-$\Lambda$-vacuum.
        This must not be papered over with a reference choice. The most the route can
        then claim is \emph{linearized matter-response} (a spin-2 perturbation) on a
        FIXED non-Einstein $\mathbb{R}\times H^3$ background -- a strictly weaker result
        that must be NAMED as such, never relabeled ``GR derived.''
  \item \textbf{If the $\eta+h$ bridge:} the thesis is wrong as stated and must be
        restated -- show why the bridge metric (not the cone curvature) is the
        gravitational field, carry the built-in $\Lambda=0$ tripwire into Phase 73, and
        (b) becomes the reference.
\end{itemize}

\paragraph{Discipline for the reopened Phase 70.} Exact over $\mathbb{Q}$; do NOT
insert any factor to force $\mathrm{Ric}\propto g$ or to force flatness (that corrupts
the very verdict it feeds); $\det_3$ remains the single source of truth; headline the
eigenvalues $\{0,-1,-1,-1\}$.

\subsection{What is verified and exact (carry-forward to 72-02)}
\label{sec:carryforward}

\begin{itemize}
  \item Peirce index map $V_1=\{0\}$, $V_0=\{1..10\}$, $V_{1/2}=\{11..26\}$
        ($L_{E_{11}}$ diagonal); cross-term $2\,\mathrm{Re}((x_2x_1)x_3)$ the unique
        $V_0\!\leftrightarrow\!V_{1/2}$ octonion channel; $\alpha$ absent from it.
        [exact over $\mathbb{Q}$]
  \item Representative non-vacuous $M$ (\texttt{bg\_delta}, \texttt{matter\_delta}
        above) with cross-term $-\tfrac{13}{315}\neq0$, all three slots, $e_4$ content.
  \item Phase-71 matterless $R$ anchors reproduced exactly (regression passes); the
        common slice point is \texttt{ROUTE1\_COMMON\_SLICE\_PT}, not the (degenerate
        off-center) center slice.
  \item $n=4$ Ricci decomposition (scalar + traceless-Ricci + Weyl) implemented and
        reconstruction-consistent ($R_{ijkl}-(\text{Scal}+E+\text{Weyl})=0$ over
        $\mathbb{Q}$ at the center); it correctly reports the cone-Hessian center as
        NON-Einstein ($S_{\mu\nu}\neq0$), which is what surfaced the obstruction.
  \item Isolation helper $\Delta R(M)=R(X_{\mathrm{bg}}+M)-R(X_{\mathrm{bg}})$ at a
        fixed $V_0$-background returns real, nondegenerate, exact-over-$\mathbb{Q}$
        values, and $\Delta R(M{=}0)=0$ (trivial limit) -- PRESERVED and unaffected by
        the backtrack; its decisive interpretation awaits the reopened Phase 70's
        metric selection (Sec~\ref{sec:phase70-task}). The matter mechanism is not what
        broke.
\end{itemize}
